\documentclass[11pt]{article}
\usepackage{fontspec}
\usepackage[french,english]{babel}
\usepackage{amsmath}
\usepackage{amssymb}
\usepackage{amsthm}
\usepackage{mathrsfs}
\usepackage{tikz-cd}
\usepackage{enumitem}
\usetikzlibrary{babel}

\newcommand{\printedpage}[1]{\marginpar{\footnotesize p.~#1}}
\newcommand{\msp}[1]{\marginpar{\footnotesize [#1]}}

\newcommand{\cM}{\mathcal{M}}
\newcommand{\cN}{\mathcal{N}}
\newcommand{\cX}{\mathcal{X}}
\newcommand{\cY}{\mathcal{Y}}
\newcommand{\cF}{\mathcal{F}}
\newcommand{\uHom}{\underline{\mathrm{Hom}}}
\newcommand{\Hom}{\mathrm{Hom}}
\newcommand{\Ens}{\mathrm{Ens}}
\newcommand{\Top}{\mathrm{Top}}
\newcommand{\Ob}{\mathrm{Ob}\,}
\newcommand{\defeq}{\overset{\text{def}}{=}}
\newcommand{\bt}{\boxtimes}
\newcommand{\tbt}{\mathbin{\widetilde{\boxtimes}}}
\newcommand{\linj}{\lhook\joinrel\longrightarrow}

\begin{document}

\printedpage{4}
% head: A. Grothendieck : Dérivateurs, Ch. XVIII, Édition des Universités de Montpellier II et Paris VII

\msp{page 57}

\[
\begin{tikzcd}[column sep=small]
 & \cM \tbt \cN \arrow[dl, "\sim" sloped] \arrow[dr, "\sim" sloped] & \\
\uHom^{!}(\cM^{\circ},\cN) \arrow[dr, "\sim" sloped] & & \uHom^{!}(\cN^{\circ},\cM) \arrow[dl, "\sim" sloped] \\
 & \underbrace{\cM \bt \cN}_{\defeq \uHom^{!!}(\cM^{\circ}\times\cN^{\circ},\Ens)} &
\end{tikzcd}
\tag*{(4.1.16)}
\]

\texttt{[where $\cM \tbt \cN$ is the category of pairs defined on the preceding page.]}

I say that $\cM \bt \cN$ \textit{is a pseudo-topos category}. It is clear, \texttt{[in]} addition, that it admits small $\varprojlim$ , which are computed argument by argument in $\uHom(\cM^{\circ}\times\cN^{\circ},\Ens)$, or at one's choice in $\uHom^{!}(\cM^{\circ},\cN)$, or $\uHom^{!}(\cN^{\circ},\cM)$. How to construct small $\varinjlim$ ? There is no question of invoking an argument using cogenerators. Besides, one must also find a small strictly generating subcategory. Evidently, there is something nontrivial to prove!

\section*{4.2\quad Muddles over $\cM \bt \cN$}

Let us return to the general situation of a pseudo-topos $\cX$, described by the pseudo-topos category $\cM$. If $\cN$ is any category, we define the category

\msp{page 58}
\[
\mathrm{Fais}(\cX,\cN) \qquad \text{or} \qquad \cF(\cX,\cN)
\]
of sheaves on $\cX$, with values in $\cN$, by
\[
\cF(\cX,\cN) \ \defeq\ \uHom^{!}(\cM^{\circ},\cN) \ .
\tag*{(4.2.1)}
\]
We thus have a fully faithful functor
\[
\cF(\cX,\cN) \ \defeq\ \uHom^{!}(\cM^{\circ},\cN) \ \linj\ \Hom^{!}(\cN^{\circ},\overline{\cM}) \ ,
\tag*{(4.2.2)}
\]
where $\overline{\cM}$ denotes the subcategory of $\cM^{\wedge}$ (\textbf{N.B.} $\cM^{\wedge}$ is not in general a $\mathcal{U}$-category never mind that) formed by the \textit{representable} functors $\cM^{\circ}\longrightarrow\Ens$ (which is a $\mathfrak{U}$-category, equivalent to $\cM$ by the canonical inclusion $\cM \linj \overline{\cM}$). Thus we have up to unique isomorphism a factorization of (4.2.2) as
\[
\cF(\cX,\cN) \ \overset{\text{fully faithful}}{\lhook\joinrel\relbar\joinrel\relbar\joinrel\longrightarrow}\ \uHom^{!}(\cN^{\circ},\cM) \ ,
\tag*{(4.2.3)}
\]
the essential image consisting of the functors
\[
\psi \ :\ \cN^{\circ} \longrightarrow \cM
\]
such that for every $x$ in $\cM$,
\[
y \longmapsto \Hom(y,\psi(x)) \ :\ \cN^{\circ} \longrightarrow \Ens
\]
is representable, i.e. of functors which
% foot: Version 17 décembre 2025
\printedpage{5}
% head: A. Grothendieck : Dérivateurs, Ch. XVIII, Édition des Universités de Montpellier II et Paris VII
\msp{page 59}
admit a \textit{left adjoint}. When $\cN$ itself is a pseudo-topos category, this condition is automatic (it follows from the hypothesis that $\psi$ commutes with small $\varprojlim$ ). Thus in this case, (4.2.3) is an equivalence of categories.

Now let $S$ ($S$ as in ``site'') be a \textit{small} subcategory of $\cM$, generating by strict epimorphisms. We know that this means that the composite functor
\[
\cM \xrightarrow{\ \sim\ } \uHom^{!}(\cM^{\circ},\Ens) \xrightarrow{\text{restriction to } S} S^{\wedge} \ (\ \defeq \uHom(S^{\circ},\Ens)\ )
\tag*{(4.2.4)}
\]
is fully faithful, and this implies more generally, that for any category $\cN$, the restriction functor
\[
\cF(\cX,\cN) \ \defeq\ \uHom^{!}(\cM^{\circ},\cN) \ \overset{\text{fully faithful}}{\lhook\joinrel\relbar\joinrel\relbar\joinrel\longrightarrow}\ \uHom(S^{\circ},\cN)
\tag*{(4.2.5)}
\]
is fully faithful. Indeed, the second member, via the fully faithful embedding
\[
\cN \ \overset{\text{fully faithful}}{\lhook\joinrel\relbar\joinrel\relbar\joinrel\longrightarrow}\ \uHom^{!}(\cN^{\circ},\Ens) \ \linj\ \cN^{\wedge} \ ,
\]
embeds into
\[
\uHom(S^{\circ},\uHom^{!}(\cN^{\circ},\Ens)) \ \simeq\ \Hom^{!}(\cN^{\circ},S^{\wedge}) \ ,
\]
and the full faithfulness of (4.2.5) is therefore equivalent to that of the composite functor
\[
\cF(\cX,\cN) \ \defeq\ \uHom^{!}(\cM^{\circ},\cN) \ \longrightarrow\ \uHom^{!}(\cN^{\circ},S^{\wedge}) \ \simeq\ \uHom(S^{\circ},\uHom^{!}(\cN^{\circ},\Ens)) \ .
\tag*{(4.2.6)}
\]

\msp{page 60}
Now we have commutativity in the square
\[
\begin{tikzcd}[column sep=large, row sep=large]
\cF(\cX,\cN) \defeq \uHom^{!}(\cM^{\circ},\cN) \arrow[r, "\text{\scriptsize(4.2.5)}"] \arrow[d, hook, "\text{\scriptsize(4.2.3)}"', "\text{\scriptsize fully faithful}"] & \uHom(S^{\circ},\cN) \arrow[d, hook, "\text{\scriptsize fully faithful}"] \\
\uHom^{!}(\cN^{\circ},\cM) \arrow[r, hook, "\text{\scriptsize(4.2.4)}", "\text{\scriptsize fully faithful}"'] & \underbrace{\uHom^{!}(\cN^{\circ},S^{\wedge})}_{\simeq\, \uHom(S^{\circ},\uHom^{!}(\cN^{\circ},\Ens))}
\end{tikzcd}
\tag*{(4.2.7)}
\]
(where the vertical arrows are equivalences, if $\cN$ itself is a pseudo-topos category). Let us note on the other hand that we have an equivalence
\[
\uHom(S^{\circ},\cN) \ \xleftarrow{\ \sim\ }\ \uHom^{!}(S^{\wedge\circ},\cN) \ \defeq\ \cF(\Top(S),\cN) \ ,
\tag*{(4.2.8)}
\]
where $\Top(S)$ is the pseudo-topos (in fact, a topos) defined by the pseudo-topos category $S^{\wedge}$. For this, let us note that we have
\[
\uHom^{!}(S^{\wedge\circ},\cN) \ \xrightarrow{\ \sim\ }\ \uHom_{!}(S^{\wedge},\cN^{\circ})^{\circ} \ ,
\]
and on the other hand the restriction functor to $S \linj S^{\wedge}$ gives
\[
\uHom_{!}(S^{\wedge},\cN^{\circ}) \ \xrightarrow{\ \sim\ }\ \uHom(S,\cN^{\circ}) \ ,
\]
whence
% foot: Version 17 décembre 2025
\printedpage{6}
% head: A. Grothendieck : Dérivateurs, Ch. XVIII, Édition des Universités de Montpellier II et Paris VII
\msp{page 61}
\[
\uHom(S^{\wedge\circ},\cN) \ \xrightarrow{\ \sim\ }\ (\uHom(S,\cN^{\circ}))^{\circ} \ \xrightarrow{\ \sim\ }\ \uHom(S^{\circ},\cN) \ ,
\]
which is none other than (4.2.8). Thus:

\textbf{Proposition 4.2.1.} \textit{Let $\cM=\cF(\cX)$ be a pseudo-topos category, equipped with a strictly generating full subcategory $S$, $\cN$ any $\mathfrak{U}$-category, then the following ``restriction to $S$'' functor is fully faithful.}
\[
\cF(\cX,\cN) \ \defeq\ \uHom^{!}(\cM^{\circ},\cN) \ \longrightarrow\ \uHom(S^{\circ},\cN) \ \xleftarrow{\ \sim\ }\ \underbrace{\uHom(S^{\wedge\circ},\cN)}_{\defeq\, \cF(\Top(S),\cN)} \ ,
\tag*{(4.2.1.1)}
\]
\textit{and it fits into the commutative square (4.2.7) of fully faithful functors.}

\textbf{Corollary 4.2.2.} \textit{Suppose that $\cN=\cF(\cY)$ is also a pseudo-topos category, equipped with a strictly generating full subcategory $T$. Then}
\[
\cF(\cX,\cN) \ \simeq\ \cF(\cY,\cM) \ \simeq\ \underbrace{\uHom^{!!}(\cM^{\circ}\times\cN^{\circ},\Ens)}_{=\,\cM\bt\cN} \ ,
\]
\textit{and the restriction functor}
\[
\begin{aligned}
\cM\bt\cN \ \defeq\ \uHom^{!!}(\cM^{\circ}\times\cN^{\circ},\Ens) \quad &\longrightarrow \quad \uHom(S^{\circ}\times T^{\circ},\Ens)\\
&\simeq \quad (S\times T)^{\wedge}\\
&\defeq \quad \cF(\Top(S\times T))
\end{aligned}
\tag*{(4.2.2.1)}
\]
\textit{is fully faithful.}

The first assertion was seen in 4.1. For the full faithfulness of (4.2.2.1), we note that this functor is interpreted, up to equivalence of categories, as the composite of the fully faithful functor (4.2.1.1), and of the functor
\msp{page 62}
analogous for $\cN$ (hence it too fully faithful)
\[
\cF(\cY,S^{\wedge}) \ \defeq\ \uHom^{!}(\cN^{\wedge},S^{\wedge}) \ \xrightarrow{\ \text{restr. to } T\ }\ \underbrace{\uHom(T^{\circ},S^{\wedge})}_{\uHom(S^{\circ}\times T^{\circ},\Ens)} \ ,
\tag*{(4.2.2.2)} % [sic] \uHom^{!}(\cN^{\wedge},S^{\wedge}) -> \uHom^{!}(\cN^{\circ},S^{\wedge}), cf. (4.2.2.3) where the same F(Y,S^) is Hom!(N°,S^)
\]
taking into account the symmetry equivalence
\[
\cF(\cY,S^{\wedge}) \ \defeq\ \uHom^{!}(\cN^{\circ},S^{\wedge}) \ \simeq\ \uHom^{!}(S^{\wedge\circ},\cN) \ .
\tag*{(4.2.2.3)}
\]

All in all, when $\cM$, $\cN$ are both pseudo-topos categories, equipped with strictly generating subcategories $S$ resp. $T$, we have a canonical fully faithful embedding
\[
\cM\bt\cN \ \overset{\rho}{\linj}\ (S\times T)^{\wedge} \ ,
\tag*{(4.2.9)}
\]
% foot: Version 17 décembre 2025
\printedpage{7}
% head: A. Grothendieck : Dérivateurs, Ch. XVIII, Édition des Universités de Montpellier II et Paris VII
which we subsequently interpret as $(f_{S}^{\cM}\times f_{T}^{\cN})_{*}$, where
\[
\left\{
\begin{aligned}
f_{S}^{\cM} \ &:\ \cX \ \longrightarrow\ \Top S\\
f_{T}^{\cN} \ &:\ \cY \ \longrightarrow\ \Top T
\end{aligned}
\right.
\tag*{(4.2.10)}
\]
are the morphisms of pseudo-topoi, given by the fully faithful embeddings (direct image functors)
\[
\left\{
\begin{aligned}
(f_{S}^{\cM})_{*} \ &:\ \cM \ \longrightarrow\ S^{\wedge}\\
(f_{T}^{\cN})_{*} \ &:\ \cN \ \longrightarrow\ T^{\wedge} \ ,
\end{aligned}
\right.
\tag*{(4.2.11)}
\]
which do indeed admit left adjoints
\msp{page 63}
\[
\left\{
\begin{aligned}
(f_{S}^{\cM})^{*} \ &:\ S^{\wedge} \ \longrightarrow\ \cM\\
(f_{T}^{\cN})^{*} \ &:\ T^{\wedge} \ \longrightarrow\ \cN \ ,
\end{aligned}
\right.
\tag*{(4.2.12)}
\]
namely the canonical extensions (by commutation with small $\varinjlim$ ) of the inclusion functors $S \linj \cM$, $T \linj \cN$.

The essential question, now, is evidently to construct a left adjoint for the functor (4.2.9). It amounts to the same thing, for $(s,t)$ in $S\times T$, to define an object
\[
s\bt t \ \in\ \Ob\cM\bt\cN \ ,
\tag*{(4.2.13)}
\]
such that we have, for $F$ varying in $\cM\bt\cN$, an isomorphism functorial in $F$,
\[
\underbrace{\rho(F)(s,t)}_{\defeq\, F(s,t)} \ \simeq\ \Hom(s\bt t,F) \ .
\]
Thus we must prove this:

\textbf{Lemma 4.2.3} (\footnote[4]{Not proved; perhaps false!}). \textit{Let $\cM$, $\cN$ be two pseudo-topos categories, and $s\in\Ob\cM$, $t\in\Ob\cN$, and let us consider on $\cM\bt\cN\defeq\uHom^{!!}(\cM^{\circ}\times\cN^{\circ},\Ens)$ the functor $F\longmapsto F(s,t)$. This functor is} representable \textit{(by an object which I denote by $s\bt t$).}

The natural idea would be to define $s\bt t$ by
\[
(s\bt t)(x,y) \ =\ s(x)\times t(y) \ =\ \Hom_{\cM}(x,s)\times\Hom_{\cN}(y,t) \ ,
\]

\msp{page 64}
which does define a functor
\[
\cM^{\circ}\times\cN^{\circ} \longrightarrow \Ens \ .
\]
Unfortunately, for one of the arguments fixed, say for $y$ fixed, this functor is \textit{not} compatible with inverse limits, e.g. it does not transform the final object of $\cM^{\circ}$ corresponding to the initial object of $\cM$) into a final object of Ens, but into $t(y)=\Hom_{\cN}(y,t)$. % [sic] "de $\cM$)" -> closing parenthesis without opening one, unbalanced as printed
% foot: Version 17 décembre 2025
\printedpage{8}
% head: A. Grothendieck : Dérivateurs, Ch. XVIII, Édition des Universités de Montpellier II et Paris VII
Nevertheless, I think that $s\bt t$ must indeed be definable as the product of two objects of $\cM\bt\cN$, depending functorially on $s$ resp. on $t$ (\footnote[5]{also false, cf. below.}). We are thus led right now, whether we like it or not, to try to define canonical functors, commuting with small $\varinjlim$ ,
\[
\cM \longrightarrow \cM\bt\cN \longleftarrow \cN
\tag*{(4.2.13)} % [sic] (4.2.13) -> (4.2.14), the label (4.2.13) is already used for s⊠t ∈ Ob M⊠N
\]
(which will be the inverse image functors for the canonical projections
\begin{tikzcd}[row sep=small, column sep=small, baseline=(\tikzcdmatrixname-2-1.base)]
 & \cX \\
\cX\times\cY \arrow[ur] \arrow[dr] & \\
 & \cY
\end{tikzcd}
) (\footnote[6]{The case where $\cM$, $\cN$ are topos-categories, shows us that these functors should be described as $x\longmapsto x\bt e_{\cN}$, $y\longmapsto e_{\cM}\bt y$. But it is not clear that we have outside the case of topoi, and even granting the existence of the $x\bt y$, the formula (valid for a topos) $x\bt y\simeq(x\bt e_{\cN})\times(e_{\cM}\bt y)$.}).

\section*{4.3\quad The operation $\cM\bt\cN$ and the accessibility axiom}

I spent the day yesterday, trying to establish the existence of the $x\bt y$ in $\cM\bt\cN$, without succeeding – at least not without making an additional accessibility hypothesis,
\msp{page 65}
about which it is becoming obvious now that it has to be made in any case. I am going all the same to state the positive of my perplexities.

\textbf{Proposition-perplexity 4.3.1.} \textit{Let $\cM$, $\cN$ be two pseudo-topos categories, whence $\cM\bt\cN$ as in 4.1. Let $S\subset\cM$, $T\subset\cN$ be two small strictly generating subcategories, and let us consider the corresponding fully faithful inclusion:}
\[
\begin{tikzcd}[column sep=large]
\underbrace{\cM\bt\cN}_{\defeq\, \uHom^{!!}(\cM^{\circ}\times\cN^{\circ},\Ens)} \arrow[r, hook, "\alpha"'] & S^{\wedge}\bt T^{\wedge} \ \simeq\ (S\times T)^{\wedge}. \arrow[l, dotted, bend right, "\varphi\,?"']
\end{tikzcd}
\tag*{(4.3.1.1)}
\]
\textit{Let us consider the following conditions:}
\begin{enumerate}[label={}]
\item[\textit{a)}] \textit{$\cM\bt\cN$ is a pseudo-topos.}
\item[\textit{b)}] \textit{For $x\in\Ob\cM$, $y\in\Ob\cN$, $x\bt y\in\Ob\cM\bt\cN$ exists, i.e. the functor (commuting with small $\varprojlim$ )}
\[
F \ \longmapsto\ F(x,y) \ , \qquad \cM\bt\cN \longrightarrow \Ens
\tag*{(4.3.1.2)}
\]
\textit{is representable.}
\item[\textit{b$'$)}] \textit{As in b), with $x\in\Ob S$, $y\in\Ob T$.}
\item[\textit{c)}] \textit{The functor $\alpha$ (which commutes with small $\varprojlim$ ) admits a left adjoint}
\[
\varphi \ :\ (S\times T)^{\wedge} \longrightarrow \cM\bt\cN \ .
\]
\end{enumerate}
\textit{We then have the implications}
\[
b \ \Longrightarrow\ b' \ \Longleftrightarrow\ c \ \Longrightarrow\ a \ .
\tag*{(4.3.1.3)}
\]
% foot: Version 17 décembre 2025

\end{document}
